\minitoc 

% ========================================================================================================
% Introduction 


% ========================================================================================================
% Systèmes et signaux à temps discret 

\section{Systèmes et signaux à temps discret}

\subsection{Généralités}

\begin{definition}[Signal]
    On définit un signal numérique comme une suite $(x_n)_{n \in \Z} $ 
    à valeurs dans $ \C$. 
    En pratique, elle sera finie et sera stockée sous forme de vecteur. 
\end{definition}

\begin{example}
    Voyons quelques exemples de signaux caractéristiques : 
    \begin{itemize}
        \item \textbf{L'échantillon unitaire}, aussi appelé Dirac : 
            \begin{equation}
                x(n) = \delta(n) = 
                    \begin{cases}
                        1 \text{ si } n = 0 \\ 
                        0 \text{ sinon } 
                    \end{cases}
            \end{equation}
        \item \textbf{Fonction de Heaviside} : 
            \begin{equation}
                x(n) = u(n) = 
                    \begin{cases}
                        1 \text{ si }n \geqslant 0 \\
                        0 \text{ sinon}  
                    \end{cases}
            \end{equation}
        \item \textbf{Sinon/Cosinus réel} : 
            \begin{equation}
                x(n) = A \cos (2 \pi f_0)n 
            \end{equation}
            Avec $A \in \R$ l'amplitude du signal. 
        \item \textbf{Exponentiel réel} : 
            \begin{equation}
                x(n) = a^n, \quad a \in \R 
            \end{equation}
        \item \textbf{Exponentiel complexe} : 
            \begin{equation}
                x(n) = c^n, \quad c \in \C
            \end{equation}
            Si $c = e^{j \pi f_0}$ (i.e $c$ appartient au cercle unité), alors : 
                $$ x(n) = \cos ( \pi n f_0) + j \sin ( \pi n f_0) $$ 
            où $j$ est la racine cubique de l'unité. 
        \item \textbf{Porte} : 
            \begin{equation}
                x(n) = 
                \begin{cases}
                    1 \text{ si } N \leqslant n \leqslant M \\ 
                    0 \text{ sinon }  
                \end{cases}
                \quad N, M \in \Z 
            \end{equation}
    \end{itemize}
\end{example}

\begin{remark}
    On peut remarquer que l'échantillon unitaire peut s'écrire en fonction de la fonction 
    de Heaviside sous la forme : 
        $$ \forall n \in \Z, \quad \delta(n) = u(n) - u(n-1) $$ 
    De même, on a : 
        $$ u(n) = \sum_{k = - \infty}^{n} \delta(k) $$ 
\end{remark}

\begin{definition}[Signal fini/périodique]
    Un signal $(x_n)_{n \in \Z}$ est dit \emph{de longueur finie} s'il 
    existe $ N_1, N_2 \in \Z$ tels que : 
        \begin{equation}
            \forall n \not \in \llbracket N_1, N_2 \rrbracket, x(n) = 0 
        \end{equation}
    On dira alors que le signal a une longueur de $N_2 - N_1 + 1$. 
    On dira aussi qu'il est périodique s'il existe $ N \in \N$ tel que : 
        \begin{equation}
            \forall n \in \Z, \quad x(n + N) = x(n) 
        \end{equation}
\end{definition}

\subsection{Opérations}

Définissons maintenant quelques opérations sur les signaux. 

\begin{definition}[Opérations]
    Soit $(x_n)_{n \in \Z}$ un signal à valeurs dans $\C$. On définit les opérations suivantes : 
    \begin{itemize} 
        \item \textbf{Décalage} : Pour $n_0 \in \Z$, on a : 
            $$ \forall n \in \Z, \quad y(n) = x(n - n_0) $$ 
        Si $n_0 > 0$ alors $y$ est en retard par rapport à $x$, sinon il est en 
        avance. 
        \item \textbf{Retournement} : $ \forall n \in \Z, y(n) = x(-n) $. 
        \item \textbf{Sous-échantillonnage} : 
            $$ y(n) = x(M \times n) $$
        où $M > 0$ est appelé \emph{facteur de sous-échantillonnage}. 
        Cette opération permet de sélectionner quelques valeurs du signal d'origine 
        périodiquement. 
        \item \textbf{Sur-échantillonnage} :
            $$ \forall n \in \Z, \quad 
                y(n) = x \left(n / N\right) 
                \quad \iff \quad 
                y(n) = 
                \begin{cases}
                    x(n/N) \text{ si } n = 0, \pm N, \pm 2N... \\ 
                    0 \text{ sinon }  
                \end{cases}
            $$
    \end{itemize}
\end{definition}

\begin{theorem}[Décomposition]
    Tout signal $(x_n)_{n \in \Z} \in \C^\Z$ peut être décomposé en somme de Dirac : 
        \begin{equation}
            \boxed{\forall n \in \Z, \quad x(n) = \sum_{k \in \Z} x(k) \delta(n - k)}
        \end{equation}
\end{theorem}

\subsection{Systèmes}

\begin{definition}[Système]
    Un système $T$ est une transformation qui modifie un signal $(x_n)_{n \in \Z}$ en 
    $T[(x_n)_{n \in \Z}]$. 

    On dira que $T$ est \emph{linéaire} si $ \forall \alpha_1, \alpha_2 \in \C$ et 
    $ \forall x_1, x_2 \in \C^\Z$, on a :
        \begin{equation}
            \forall n \in \Z, \quad T[ \alpha_1, x_1(n) + \alpha_2 x_2(n)] = \alpha_1 T[x_1(n)] + \alpha_2 T[x_2(n)] 
        \end{equation}
    On appelle cela le \emph{principe de superposition}. 
\end{definition}

\begin{prop}[Systèmes à temps discret]
    Soient $x,y \in \C^\Z$ deux signaux et $T$ un système quelconque. 
    On définit les propriétés suivantes pour $T$ : 
    \begin{itemize}
        \item \textbf{Causalité} : Si pour tout $n_0 \in \Z$, $T[x(n_0)]$ 
        ne dépend que de $x(n)$ et/ou $T[x(n)]$ avec $n \leqslant n_0$ . 
        \item \textbf{Invariance par décalage} : si 
            \begin{equation}
                \forall n_0 \in \Z, \forall n \in \Z, \quad T[x(n - n_0)] = T[x(n)] 
            \end{equation}
        Dans le cas contraire, on dira que $T$ est variant par décalage. 
        \item \textbf{Stabillité} : Si pour toute entrée bornée, la sortie est elle-même bornée. 
        Plus formellement, il existe $A,B \in \R_+$ tels que : 
            \begin{equation}
                \forall n \in \Z, \quad |x(n)| < A \Longrightarrow T[x(n)] < B 
            \end{equation}
    \end{itemize}
\end{prop}

On appellera SLIT un \textbf{S}ystème \textbf{L}inéaire \textbf{I}nvariant par \textbf{T}ranslation. 

\begin{definition}[Système physiquement réalisable]
    On dira qu'un système $T$ est physiquement réalisable, s'il est à la fois \emph{stable} 
    et \emph{causal}. 
\end{definition}

\begin{definition}[Réponse impultionnelle]
    On appelle \emph{réponse impultionnelle} le signal $(h_n)_{n \in \Z}$ pour un système 
    $T$ définit par : 
        \begin{equation}
            \forall n \in \Z, \quad h(n) = T[ \delta(n)]
        \end{equation}
\end{definition}

Les systèmes linéaires disposent de très bonnes propriétés et leur stabilité est facilement 
caractérisable. 

\begin{proposition}
    Soient $T$ un système linéaire de réponse impulsionnelle $(h_n)_{n \in \Z}$, alors 
    $T$ est \emph{stable} si 
    \begin{equation}
        \exists A \in \R_+ \quad \forall n \in \Z, \sum_{n \in \Z} |h(n)| \leqslant A 
    \end{equation}
\end{proposition}


% ========================================================================================================
% Convolution et corrélation

\section{Convolution et corrélation}

Pour tout signal $(x_n)_{n \in \Z}$, on sait que : 
    $$ \forall n \in \Z, \quad x(n) = \sum_{k \in \Z} x(k) \delta(n - k) $$ 
Soit $T$ un SLIT. On a alors : 
    \begin{align*}
        \forall n \in \Z, y(n) &= T[x(n)] = T\left[ \sum_{k \in Z} x(k) \delta(n - k)\right] \\ 
        &= \sum_{k \in \Z} x(k) T[ \delta(n - k)]
    \end{align*}
Or $T[ \delta(n)] = h(n)$ donc posons : 
    $$ \forall n \in \Z, \forall k \in \Z, \quad T[ \delta(n - k)] = h_k (n) $$ 
Or puisque $T$ est invariant par translation, on a : 
    $$ T[ \delta(n - k)] = h_k(n) = h(n - k) $$
Donc pour un SLIT : 
    $$ y(n) = \sum_{k \in \Z} x(k) h(n - k) $$
La réponse impultionnelle contient donc énormément d'information sur le système $T$ considéré. 
Cette formule nous donne aussi la définition de la convolution temporelle discrète. 

\subsection{Convolution à temps discret}

Dans cette section, on suppose que tous les signaux consédérés ont des propriétés pertinentes telles 
que la finitude. 

\begin{definition}[Convolution]
    Soient $(x_n)_{n \in \Z}$ et $ (y_n)_{n \in \Z}$ deux signaux. On définit 
    la \emph{convolution} entre $x$ et $y$ comme l'opération : 
        \begin{equation}
            \boxed{\forall n \in \Z, \quad (x * y)(n) = \sum_{k \in \Z} x(k) y(n - k)}
        \end{equation}
\end{definition}

\begin{prop}[Convolution et opérations]
    Soient $x, y, z \in \C^\Z$ des signaux. 
    L'opérateur convolution possède les propriétés suivantes : 
    \begin{enumerate}[label=\roman*)]
        \item \emph{Commutativité} : $ x * y = y * x $. 
        \item \emph{Associativité} : $ (x * y) * z = x * (y * z)$. 
        \item \emph{Distributivité} : $ x * (y + z) = x * y + x * z$. 
    \end{enumerate}
\end{prop}

\begin{prop}[Convolution et Dirac]
    Soit $ (x_n)_{n \in \Z}$ un signal. On a les propriétés suivantes vis-à-vis de l'échantillon 
    unitaire : 
    \begin{enumerate}[label=\roman*)]
        \item \textbf{Invariance} : 
            \begin{equation}
                \forall n \in \Z, \quad (x * \delta)(n) = \sum_{k \in \Z} x(k) \delta( n -k ) = x(n) 
            \end{equation}
        \item \textbf{Translation} : pour tout $n, n_0 \in \Z$, 
            \begin{align*}
                x(n) * \delta(n - n_0) &= x(n - n_0) \\ 
                \delta(n - n_0) * \delta(n - n_1) &= \delta(n - n_0 - n_1)
            \end{align*}
    \end{enumerate}
\end{prop}

\subsection{Cas des signaux de longueur finie}

\begin{remark}
    Soient $x, y \in \C^\Z$ deux signaux de longueur finie respective $N_x, N_y \in \N$. 
    Par définition de la convolution, pour tout $n \in \Z$, on a : 
        $$ z(n) = (x * y)(n) = \sum_{k \in \Z} x(k) y(n - k) $$
    Alors $(z_n)_{n \in \Z}$ est de longueur, au maximum $N_x + N_y - 1$.
\end{remark}

On cherche ici à définir une convolution dite \emph{circulaire} dont la longueur du résultat 
est la même que celle des opérandes. 

\begin{definition}[Extension périodique]
    Soit $(x_n)_{n \in \Z}$ un signal de longueur finie $N$. On définit son extension 
    périodique comme le signal $({x_N}_n)_{n \in \Z}$ définit par : 
        \begin{equation} 
            \forall n \in \Z, \quad {x}_N(n) = \sum_{k=-\infty}^{+\infty} x(n - kN) 
        \end{equation}
    Ou de manière équivalente : 
        \begin{equation} 
            \forall n \in \Z, \quad {x}_N(n) = x(n \mod N) 
        \end{equation}
    Ce signal est périodique de période $N$. 
\end{definition}

\begin{definition}[Convolution circulaire]
    Soient $x, y \in \C^\Z$ deux signaux de longueur finie $N \in \N$. Soit $x_N \in \C^\Z$ 
    l'extension périodique de $x$. On définit la convolution circulaire de $x$ et $y$ comme : 
        \begin{equation}
            \boxed{\forall n \in \Z, \quad (x \circ y)(x) = \sum_{k=0}^{N-1} y(k) x_N(n - k)}
        \end{equation}
\end{definition}

% \begin{center}
%     \begin{tikzpicture}[scale=2]
%         % Styles
%         \tikzstyle{val_ext} = [circle, fill=blue!10, draw=blue!80, inner sep=1.5pt, font=\footnotesize]
%         \tikzstyle{val_int} = [circle, fill=red!10, draw=red!80, inner sep=1.5pt, font=\footnotesize]

%         % Disque Extérieur (Signal x)
%         \fill[blue!5, draw=blue!30, thick] (0,0) circle (1.2cm);
%         \node[blue!70, above] at (0, 1.25cm) {Disque fixe $x(k)$};

%         % Disque Intérieur (Signal h retourné et décalé)
%         \fill[red!5, draw=red!30, thick] (0,0) circle (0.75cm);
%         \node[red!70, below] at (0, -1.35cm) {Disque mobile $h((n-k)_N)$};

%         % Échantillons et graduations (N=4)
%         % Pour n=0 : x(0)*h(0) + x(1)*h(3) + x(2)*h(2) + x(3)*h(1)
        
%         % k = 0
%         \draw[dashed, gray!50] (0,0) -- (0:1.2cm);
%         \node[val_ext, blue] at (0:1cm) {$x(0)$};
%         \node[val_int, red] at (0:0.55cm) {$h(0)$};

%         % k = 1
%         \draw[dashed, gray!50] (0,0) -- (90:1.2cm);
%         \node[val_ext, blue] at (90:1cm) {$x(3)$};
%         \node[val_int, red] at (90:0.55cm) {$h(1)$};

%         % k = 2
%         \draw[dashed, gray!50] (0,0) -- (180:1.2cm);
%         \node[val_ext, blue] at (180:1cm) {$x(2)$};
%         \node[val_int, red] at (180:0.55cm) {$h(2)$};

%         % k = 3
%         \draw[dashed, gray!50] (0,0) -- (270:1.2cm);
%         \node[val_ext, blue] at (270:1cm) {$x(1)$};
%         \node[val_int, red] at (270:0.55cm) {$h(3)$};

%         % Flèche de rotation pour le disque intérieur
%         \draw[-{Stealth}, red!60, thick] (40:0.4cm) arc (40:320:0.4cm);
%         \node[red!60] at (0,0) {\textbf{+ n}};
        
%     \end{tikzpicture}
%     \end{center}

\subsection{Corrélation}

On souhaite maintenant définir une opération permettant de déterminer si deux signaux sont 
plus ou moins proches. 

\begin{definition}[Fonction d'intercorrélation]
    Soient $x,y \in \C^\Z$. On définit la corrélation entre $x$ et $y$ comme le signal : 
        \begin{equation}
            \forall n \in \Z, \quad C_{xy}(n) = \sum_{k \in \Z} x(k) y(k - n) 
        \end{equation}
    Si $x = y$ alors $ C_{xx}$ s'appellera la \emph{fonction d'autocorrélation}. 
\end{definition}

\begin{prop}[Corrélation]
    Soient $x,y \in \C^\Z$. La corrélation admet les propriétés suivantes : 
    \begin{enumerate}[label=\roman*)]
        \item Lien avec la convolution : 
            \begin{equation}
                \forall n \in \Z, \quad C_{xy}(n) = x(n) * y(-n)
            \end{equation}
        \item Stabilité par inversion : $ \forall n \in \Z, C_{xy}(n) = C_{xy}(-n)$.
    \end{enumerate}
\end{prop}

\begin{proposition}
    On peut aussi définir deux autres formes de fonction de corrélation pour tous signaux $x,y \in \C^\Z$ 
    de longueur finie $N$ :
    \begin{itemize}
        \item \textbf{Forme biaisée} : 
            \begin{equation}
                \forall n \in \Z, \quad C_{xy}(n) = \frac{1}{N} \sum_{k=0}^{N - n -1} x(k) y(k - n)
            \end{equation}
        \item \textbf{Forme non biaisée} : 
            \begin{equation}
                \forall n \in \Z, \quad C_{xy}(n) = \frac{1}{N - |n|} \sum_{k=0}^{N - n -1} x(k) y(k - n)
            \end{equation}
    \end{itemize}

\end{proposition}

% ========================================================================================================
% Transformée de Fourrier

\section{Transformée de Fourrier}

\subsection{Définition et Propriétés}

\begin{definition}[Transformée de Fourrier]
    On définit la transformée de Fourrier $X$ comme l'application \emph{continue} telle que : 
        $$ \boxed{X : f \longmapsto 
            \sum_{n \in \Z} x(n) e^{-2j \pi n f }} 
        $$
    On a $ X(f) \in \C$ donc, en particulier, $X(f) = |X(f)| e^{j \text{Arg} (X(f))}$. Elle est 
    aussi prédiodique de période $1$, i.e $X(f + 1) = X(f)$.

    $X(f)$ indique la "quantité" de fréquence $f$ présente dans le signal $x$ sur $ \Z$. 
    Elle donne donc des informations fréquentielles sur le signal $x(n)$. 
\end{definition}

\begin{proposition}
    Bien qu'elle ne soit pas bijective, on peut aussi définir la \emph{transformée de Fourier inverse} : 
    $$
        x(n) = \int_{-1/2}^{1/2} X(f) e^{2 j n \pi f} df 
    $$
    Pour un signal et sa transformée de Fourier, on notera $ x(x) \overset{F}{\longleftrightarrow} X(f)$. 
\end{proposition}

\begin{prop}[Transformée de Fourier]
    Soient $x,y \in \C^\Z$ deux signaux. On a les propriétés suivantes : 
    \begin{enumerate}[label=\roman*)]
        \item \emph{Linéairité} : Pour tout $ a,b \in \C$, si : 
            $$ x(n) \overset{F}{\longleftrightarrow} X(f) \quad \text{et} \quad y(n) \overset{F}{\longleftrightarrow} Y(f) $$ 
        alors : 
            $$ 
                a x(x) + b y(n) \overset{F}{\longleftrightarrow} a X(f) + b Y(f) 
            $$ 
        \item \emph{Parité} : 
            $$ 
                x(n) \overset{F}{\longleftrightarrow} X(f) \Longrightarrow x(-n) \overset{F}{\longleftrightarrow} X(-f) 
            $$
        \item Si $ x(x) \overset{F}{\longleftrightarrow} X(f)$ et $ y(n) \overset{F}{\longleftrightarrow} Y(f)$, 
        alors : 
            \begin{itemize}
                \item $x(n - n_0) \overset{F}{\longleftrightarrow} X(f) e^{-2 j \pi n_0 f}$. 
                \item $x(an) \overset{F}{\longleftrightarrow} X(f/a)$ si $a \not = 0$ (dilatation). 
                \item $ x(n) e^{2 j \pi n f_0} \overset{F}{\longleftrightarrow} X(f - f_0) $
                \item $ x(n) y(n) \overset{F}{\longleftrightarrow} X(f) * Y(f) $
                \item $ x(n) * y(n) \overset{F}{\longleftrightarrow} X(f) Y(f)$ 
            \end{itemize}
    \end{enumerate}
\end{prop}

\begin{proposition}
    On retiendra les quelques transformées de Fourier usuelles : 
    \begin{center}
        \renewcommand{\arraystretch}{2.2} % Augmente l'espace vertical pour les fractions
        \begin{tabular}{lc} % Suppression des lignes verticales pour plus de clarté
            \toprule
            \textbf{Signal temporel} $x(n)$ & \textbf{Transformée de Fourier} $X(f)$ \\ 
            \midrule
            $e^{2 j \pi f_0 n}$ & $\delta(f - f_0)$ \\ 
            $\cos(2 \pi f_0 n)$ & $\displaystyle \frac{1}{2} [ \delta(f - f_0) + \delta(f + f_0)]$ \\ 
            $\sin(2 \pi f_0 n)$ & $\displaystyle \frac{1}{2j} [ \delta(f - f_0) - \delta(f + f_0)]$ \\ 
            \bottomrule
        \end{tabular}
    \end{center}
\end{proposition}

\subsection{Transformée de Fourier à temps discret (TFD)}

Dans cette sous-section, nous ne considérons plus que des signaux 
finis $(x_0, \dots, x_{n-1})$ de longueur $ n \in \N$. 

\begin{definition}[Transformée de Fourier discrète]
    De la même façon que précédement, on définira la transformée de Fourier discrète 
    d'un signal $(x_0, \dots, x_{N-1}), N \in \N$ par : 
        $$ 
            \boxed{X(m) = \sum_{n=0}^{N-1} x(n) e^{-2 j \pi  n m / N}, \quad m \in \llbracket 0, N-1 \rrbracket.}
        $$ 
\end{definition}

\begin{remark}
    Ainsi, si $ x(n) \overset{F}{\longleftrightarrow} X(f)$, alors 
    sa transformée de Fourier discrète sera $X(m) = X(f)$ pour $ f \in \{m/N \mid m \in \llbracket 0, N-1 \rrbracket\}$. 
    $X(m)$ sera aussi périodique de période $ N$. 
\end{remark}

\begin{proposition}
    De même, la transformée de Fourier discrète inverse sera : 
        $$ x(n) = \frac{1}{N} \sum_{m=0}^{N-1} X(m) e^{2 j \pi n m / N}. $$ 
    Elle héritera des mêmes propriétés que la précédente avec la convolution circulaire. 
\end{proposition}


% ========================================================================================================
% Filtres sélectifs 

\section{Filtres sélectifs}





